\suadSeite
\subsection{Überblick und Architektur}

Das Frontend der Anwendung "`EduShelf"' wurde als moderne Single Page Application (SPA) konzipiert und implementiert. Im Gegensatz zu klassischen Webanwendungen, bei denen bei jeder Interaktion eine neue HTML-Seite vom Server geladen wird, lädt eine SPA initial ein einziges HTML-Gerüst. Die Inhalte werden anschließend dynamisch durch Client-Seitiges JavaScript nachgeladen und gerendert. Dies führt zu einer flüssigeren Benutzererfahrung, die der einer nativen Desktop-Anwendung ähnelt.

Die Architektur folgt einem modularen Komponenten-Ansatz, wobei die Quelldateien im Verzeichnis \texttt{src} klar strukturiert sind:
\begin{itemize}
    \item \textbf{components}: Beinhaltet die wiederverwendbaren UI-Komponenten (z. B. \texttt{Chat.jsx}, \texttt{Quiz.jsx}).
    \item \textbf{services}: Enthält die Logik für die Kommunikation mit dem Backend (z. B. \texttt{api.js}).
    \item \textbf{assets}: Speichert statische Ressourcen wie Bilder und Stylesheets.
\end{itemize}

\subsection{Technologie-Stack}

Für die Umsetzung des Frontends wurden folgende Technologien gewählt, um Wartbarkeit, Performance und Skalierbarkeit zu gewährleisten:

\begin{itemize}
    \item \textbf{React}: Eine von Facebook entwickelte JavaScript-Bibliothek zur Erstellung von Benutzeroberflächen. Durch das Konzept des "`Virtual DOM"' werden DOM-Manipulationen minimiert, was die Rendering-Performance signifikant erhöht. React ermöglicht zudem eine strikte Trennung von Logik und Darstellung durch Komponenten.
    \item \textbf{Vite}: Ein modernes Build-Tool, das als Ersatz für Webpack dient. Es nutzt native ES-Module im Browser für extrem kurze Startzeiten des Entwicklungsservers und bietet optimierte Builds für die Produktion.
    \item \textbf{React Router}: Diese Bibliothek übernimmt das client-seitige Routing. Sie ermöglicht die Navigation zwischen verschiedenen Ansichten (z. B. \texttt{/files}, \texttt{/chat}), ohne dass die Seite neu geladen werden muss.
    \item \textbf{Axios}: Ein Promise-basierter HTTP-Client, der für die asynchrone Kommunikation mit der REST-API des Backends verwendet wird. Er vereinfacht das Senden von Anfragen und das Abfangen von Antworten (Interceptors).
\end{itemize}

\subsection{Benutzeroberfläche und UX}

Das Design von EduShelf legt Wert auf Intuitivität und Responsivität. Die Anwendung passt sich dank CSS-Media-Queries dynamisch an verschiedene Bildschirmgrößen an.

Ein zentrales Feature ist die Unterstützung von Themes (Hell/Dunkel). Dies wird technisch über ein \texttt{data-theme}-Attribut am \texttt{document.documentElement} realisiert. Der Zustand wird im \texttt{localStorage} des Browsers persistiert, sodass die Einstellung bei einem erneuten Besuch erhalten bleibt:

\begin{lstlisting}[language=JavaScript, caption={Theme-Initialisierung in App.jsx}]
useEffect(() => {
  const currentTheme = localStorage.getItem('theme');
  if (currentTheme === 'light') {
    document.documentElement.setAttribute('data-theme', 'light');
  } else {
    document.documentElement.removeAttribute('data-theme');
  }
}, []);
\end{lstlisting}

\subsection{Kernkomponenten und Implementierung}

Die Anwendung gliedert sich in mehrere funktionale Hauptbereiche, die als React-Komponenten implementiert sind.

\subsubsection{Authentifizierung und Autorisierung}
Die Sicherheit der Anwendung wird durch eine Token-basierte Authentifizierung gewährleistet. Die Komponenten \texttt{Login.jsx} und \texttt{Register.jsx} ermöglichen den Zugang. Der aktuelle Benutzerstatus wird zentral in der \texttt{App.jsx} verwaltet. Beim Start der Anwendung wird die Route \texttt{/Users/me} aufgerufen, um eine bestehende Session zu validieren.
Zusätzlich werden rollenbasierte Zugriffe (RBAC) umgesetzt. Beispielsweise ist das \texttt{AdminPanel} nur für Benutzer erreichbar, die die Rolle "`Admin"' besitzen. Versucht ein nicht-autorisierter Nutzer, auf geschützte Routen zuzugreifen, wird er automatisch weitergeleitet.

\subsubsection{Dateiverwaltung (Files)}
Die Komponente \texttt{Files.jsx} stellt die Schnittstelle zum Dokumentenmanagement dar. Dateien können per Upload-Dialog zum Backend übertragen werden. Eine Besonderheit ist die Integration des \texttt{DocxViewer}, der eine Vorschau von Word-Dokumenten direkt im Browser ermöglicht, ohne dass diese heruntergeladen werden müssen.

\subsubsection{Lernmodule: Quiz und Flashcards}
Ein Kernaspekt von EduShelf ist das interaktive Lernen.
\begin{itemize}
    \item \textbf{Quiz}: Ermöglicht das Abfragen von Wissen in einem strukturierten Format. Die Logik für die Fragenabfolge und Auswertung ist in \texttt{Quiz.jsx} implementiert.
    \item \textbf{Flashcards}: Implementiert ein Karteikarten-System. Die Komponente \texttt{Flashcards.jsx} verwaltet den Status der Karten (Vorderseite/Rückseite) und den Lernfortschritt.
\end{itemize}

\subsubsection{Kommunikation (Chat)}
Die \texttt{Chat}-Komponente ermöglicht den Austausch von Nachrichten. Das Layout ist zweigeteilt in eine Nachrichtenliste und ein Eingabefeld. Die Nachrichten werden dynamisch gerendert, wobei zwischen eigenen Nachrichten und Empfangenen visuell unterschieden wird.

\subsection{Datenfluss und State Management}

Der Datenfluss in React ist unidirektional (von Eltern- zu Kind-Komponenten). Für den lokalen Zustand der Komponenten (z. B. Formulareingaben, Sichtbarkeit von Modals) wird der \texttt{useState}-Hook verwendet.
Seiteneffekte, wie das Laden von Daten beim Komponentenstart, werden mit dem \texttt{useEffect}-Hook gesteuert.

Die Kommunikation mit dem Server ist in der Datei \texttt{src/services/api.js} zentralisiert. Dies fördert die Wartbarkeit, da API-Endpunkte und Konfigurationen (wie Basis-URL oder Header) nur an einer Stelle definiert sind. Fehlerbehandlung (z. B. bei Verbindungsabbruch oder 401 Unauthorized Fehlern) kann so konsistent für die gesamte Anwendung umgesetzt werden.
